\section*{Sets and Indices}

\[
W=\{1,2\}
\]
set of weeks, indexed by \(w\).

\[
M=\{A,B\}
\]
set of microcomputer models.

\[
P=\{I,II\}
\]
set of processes.

\section*{Parameters}

All numerical values are read from \texttt{table\_1.csv} and \texttt{general\_parameters.csv}.

For process times per unit:
\[
a_I=\texttt{Model\_A\_hours\_per\_unit for Process I},\qquad
a_{II}=\texttt{Model\_A\_hours\_per\_unit for Process II},
\]
\[
b_I=\texttt{Model\_B\_hours\_per\_unit for Process I},\qquad
b_{II}=\texttt{Model\_B\_hours\_per\_unit for Process II}.
\]

For unit sales profit:
\[
\pi_A=\texttt{Profit\_per\_unit for Model A},\qquad
\pi_B=\texttt{Profit\_per\_unit for Model B}.
\]

Regular weekly process capacities:
\[
H_I=\texttt{process\_I\_hours},\qquad H_{II}=\texttt{process\_II\_hours}.
\]

Maximum overtime per week:
\[
\overline{O}_I=\texttt{max\_overtime\_I\_per\_week},\qquad
\overline{O}_{II}=\texttt{max\_overtime\_II\_per\_week}.
\]

Overtime cost rates:
\[
c_I=\texttt{overtime\_premium\_I},\qquad
c_{II}=\texttt{overtime\_premium\_II}.
\]

Week-1 fatigue thresholds:
\[
T_I=\texttt{week1\_overtime\_fatigue\_threshold\_I},\qquad
T_{II}=\texttt{week1\_overtime\_fatigue\_threshold\_II}.
\]

Week-2 recovery losses in regular hours:
\[
L_I=\texttt{week2\_recovery\_loss\_I},\qquad
L_{II}=\texttt{week2\_recovery\_loss\_II}.
\]

Weekly minimum delivery requirements:
\[
d_{A1}=\texttt{week1\_min\_model\_A\_units},\quad
d_{A2}=\texttt{week2\_min\_model\_A\_units},
\]
\[
d_{B1}=\texttt{week1\_min\_model\_B\_units},\quad
d_{B2}=\texttt{week2\_min\_model\_B\_units}.
\]

Two-week minimum production requirements:
\[
D_A=\texttt{min\_model\_A\_units},\qquad
D_B=\texttt{min\_model\_B\_units}.
\]

Minimum average utilization fractions of regular hours:
\[
u_I=\texttt{min\_avg\_utilization\_I},\qquad
u_{II}=\texttt{min\_avg\_utilization\_II}.
\]

Minimum required net two-week profit:
\[
\Pi^{\min}=\texttt{min\_weekly\_profit}.
\]

\section*{Decision Variables}

Production quantities:
\[
x_{Aw}\ge 0 \quad \text{(code: \texttt{xA[w]})}, \qquad
x_{Bw}\ge 0 \quad \text{(code: \texttt{xB[w]})}
\]
for \(w\in W\).

Overtime hours:
\[
o_{Iw}\ge 0 \quad \text{(code: \texttt{otI[w]})}, \qquad
o_{IIw}\ge 0 \quad \text{(code: \texttt{otII[w]})}
\]
for \(w\in W\).

Recovery activation indicators:
\[
r_I\in\{0,1\}\quad \text{(code: \texttt{recovery\_I})},\qquad
r_{II}\in\{0,1\}\quad \text{(code: \texttt{recovery\_II})}.
\]

\section*{Objective}

Maximize net two-week profit after overtime spending:
\[
\max \;
\pi_A(x_{A1}+x_{A2})+\pi_B(x_{B1}+x_{B2})
-c_I(o_{I1}+o_{I2})-c_{II}(o_{II1}+o_{II2}).
\]

\section*{Constraints}

Week 1 process capacities:
\[
a_I x_{A1}+b_I x_{B1}\le H_I+o_{I1},
\]
\[
a_{II} x_{A1}+b_{II} x_{B1}\le H_{II}+o_{II1}.
\]

Week 2 process capacities with recovery loss:
\[
a_I x_{A2}+b_I x_{B2}\le H_I-L_I r_I+o_{I2},
\]
\[
a_{II} x_{A2}+b_{II} x_{B2}\le H_{II}-L_{II} r_{II}+o_{II2}.
\]

Weekly overtime limits:
\[
o_{Iw}\le \overline{O}_I \qquad \forall w\in W,
\]
\[
o_{IIw}\le \overline{O}_{II} \qquad \forall w\in W.
\]

Week-1 overtime/recovery linking:
\[
o_{I1}\le T_I+\overline{O}_I\, r_I,
\]
\[
o_{II1}\le T_{II}+\overline{O}_{II}\, r_{II}.
\]

Weekly minimum deliveries:
\[
x_{A1}\ge d_{A1},\qquad x_{A2}\ge d_{A2},
\]
\[
x_{B1}\ge d_{B1},\qquad x_{B2}\ge d_{B2}.
\]

Two-week minimum production:
\[
x_{A1}+x_{A2}\ge D_A,
\]
\[
x_{B1}+x_{B2}\ge D_B.
\]

Minimum average utilization of regular process hours over two weeks:
\[
a_I x_{A1}+b_I x_{B1}+a_I x_{A2}+b_I x_{B2}\ge 2H_Iu_I,
\]
\[
a_{II} x_{A1}+b_{II} x_{B1}+a_{II} x_{A2}+b_{II} x_{B2}\ge 2H_{II}u_{II}.
\]

Minimum net two-week profit:
\[
\pi_A(x_{A1}+x_{A2})+\pi_B(x_{B1}+x_{B2})
-c_I(o_{I1}+o_{I2})-c_{II}(o_{II1}+o_{II2})
\ge \Pi^{\min}.
\]

\section*{Notes on Implementation}

The implemented model is a mixed-integer linear program with continuous production and overtime variables and binary recovery indicators. It is solved directly with the optimizer invoked in \texttt{current\_heuristic.py}.

The recovery logic is modeled exactly as in the code: if week-1 overtime for a process exceeds its threshold, then the corresponding binary \(r_p\) must take value \(1\), which reduces week-2 regular capacity by the fixed loss \(L_p\). However, the converse is not enforced: the formulation does not require \(r_p=0\) when week-1 overtime is at or below the threshold, so the model may choose \(r_p=1\) voluntarily, though this only tightens week-2 capacity.

The parameter named \texttt{min\_weekly\_profit} is used in the code as a lower bound on total net profit across both weeks, not as a separate per-week profit requirement.
